\documentclass[../../main.tex]{subfiles}

\begin{document}

\section{From the electron-phonon interaction to the BCS ground state}
\label{sec:froehlich-to-bcs}

\subsection{Overview} \label{subsec:froehlich-to-bcs_overview}

The road from deriving the Hamiltonian of the electron--phonon interaction to the BCS
Hamiltonian, the physics it contains and finally its ground state is quite long and
necessitates many calculations. Hence, this note only gives a concise summary and then
refers to the literature for more information. I mostly follow \emph{Solid State Theory 2}
by Gerd Czycholl; equation numbers of the form (C\,5.7) refer to that book. The
Hamiltonian used as a starting point is derived by Czycholl in Sec.~7.1 of his first book
and also by Bruus and Flensberg in Ch.~3 of \emph{Many-Body Quantum Theory in Condensed
Matter Physics}. If you are simply interested in the BCS ground state, you can directly
skip to \cref{subsubsec:bcs-ground-state_overview}. If you want the details, refer to \cref{subsec:froehlich-to-bcs_details}.

\subsubsection{Attractive electron--electron interaction through the electron-phonon
mechanism}
\label{subsubsec:attractive-interaction_overview}

\paragraph{Intuitive picture/motivation.}
An electron $e_1$ moving through the lattice distorts it, i.e.\ it drags the positive ions
towards itself and leaves behind a region of enhanced positive charge. A second electron
$e_2$ then feels an attraction to $e_1$ through this phonon excitation caused by $e_1$.
The result is an effective electron--electron interaction that is attractive.

\begin{figure}[ht]
\centering
\begin{tikzpicture}[
  x=1cm, y=1cm, line width=0.7pt,
  eqion/.style={circle, draw=gray!55, minimum size=7.5mm, inner sep=0pt},
  dsion/.style={circle, draw=teal!70!black, fill=teal!12, line width=0.8pt,
                minimum size=7.5mm, inner sep=0pt},
  elec/.style={circle, fill=black, minimum size=3pt, inner sep=0pt},
  move/.style={-{Stealth[length=5pt]}},
]
% ---------------- panel I: e_1 drags the ions along
\begin{scope}
  \node[anchor=west] at (-1.15,2.5) {I.)};
  \foreach \x/\y in {0/0, 1.7/0, 0/1.7, 1.7/1.7} {
    \node[eqion] at (\x,\y) {};
    \node[dsion] at ($ (\x,\y) + (0.17,-0.17) $) {};
  }
  \node[elec] (e1) at (0.85,0.85) {};
  \node[anchor=south west, inner sep=2pt] at (0.9,0.85) {$e_1$};
  \draw[move] (0.7,0.78) -- (-0.95,0.28);
\end{scope}
% ---------------- panel II: the distortion outlives e_1 and attracts e_2
\begin{scope}[xshift=5.6cm]
  \node[anchor=west] at (-1.15,2.5) {II.)};
  \foreach \x/\y in {0/0, 1.7/0, 0/1.7, 1.7/1.7} {
    \node[eqion] at (\x,\y) {};
    \node[dsion] at ($ (\x,\y) + (0.17,-0.17) $) {};
  }
  \node[circle, fill=red!12, minimum size=10mm, inner sep=0pt] at (0.85,0.85)
       {\textcolor{red!70!black}{$+$}};
  \node[elec] (e2) at (3.1,1.35) {};
  \node[anchor=south west, inner sep=2pt] at (3.15,1.35) {$e_2$};
  \draw[move] (2.95,1.28) -- (1.75,1.05);
\end{scope}
\end{tikzpicture}
\caption{The lattice distortion caused by $e_1$ (I.) outlives its passage and
attracts a second electron $e_2$ (II.).}
\label{fig:effective_phonon_interaction}
\end{figure}

\paragraph{Starting point.}
The starting point for all our derivations is the Fr\"ohlich Hamiltonian, derived in
Czycholl Band~1, Sec.~7.1 (7.13)--(7.15) or Bruus and Flensberg, Sec.~3.8, (C\,5.7):
\begin{equation}
H = \sum_{\vec k,\sigma}\varepsilon_{\vec k}\,c^\dagger_{\vec k,\sigma}c_{\vec k,\sigma}
  + \sum_{\vec q}\hbar\omega_{\vec q}\,b^\dagger_{\vec q}b_{\vec q}
  + \sum_{\vec k,\vec q,\sigma}\Big[
      M(\vec q)\,c^\dagger_{\vec k+\vec q,\sigma}c_{\vec k,\sigma}b_{\vec q}
    + M(-\vec q)\,c^\dagger_{\vec k-\vec q,\sigma}c_{\vec k,\sigma}b^\dagger_{\vec q}
    \Big],
\label{eq:froehlich}
\end{equation}
with $\omega_{-\vec q}=\omega_{\vec q}$ and $M(-\vec q)=M(\vec q)$.

\paragraph{Canonical transformation.}
Use a canonical transformation\footnote{See \cref{sec:similarity-transformations} for more general information on similarity transformations.} $H\rightsquigarrow H_T=e^{-iS}He^{iS}$ to find a purely
electronic Hamiltonian that describes this interaction, i.e.\ to get rid of processes that
change the number of phonons. At one point one simply restricts the Hamiltonian not to mix
states with different phonon configurations.

The canonical transformation yields $H_T\cong H_0+H_{\mathrm{int}}$ with the interaction
(C\,5.25)
\begin{equation}
H_{\mathrm{int}} = -\frac12\sum_{\sigma,\sigma'}\sum_{\vec k,\vec k',\vec q}
  V_{\vec k,\vec k',\vec q}\;
  c^\dagger_{\vec k+\vec q,\sigma}\,c^\dagger_{\vec k'-\vec q,\sigma'}\,
  c_{\vec k',\sigma'}\,c_{\vec k,\sigma},
\label{eq:h-int}
\end{equation}

which is attractive if
$|\varepsilon_{\vec k}-\varepsilon_{\vec k-\vec q}|\le\hbar\omega_{\vec q}\le\hbar\omega_D$
in $V_{\vec k,\vec k',\vec q}$, where $\omega_D$ is the Debye frequency, and can be
represented as in \cref{fig:vertex}.

\begin{figure}[ht]
\centering
\begin{tikzpicture}[
  x=1cm, y=1cm, line width=0.7pt,
  ferm/.style={decoration={markings,
               mark=at position 0.58 with {\arrow{Stealth[length=5pt]}}},
               postaction={decorate}},
]
  \coordinate (L) at (-1.05,0);
  \coordinate (R) at ( 1.05,0);
  \draw[ferm] (-2.75,-1.55) -- (L);
  \draw[ferm] (L) -- (-2.75, 1.55);
  \draw[ferm] ( 2.75,-1.55) -- (R);
  \draw[ferm] (R) -- ( 2.75, 1.55);
  \draw[decorate,
        decoration={coil, aspect=0, amplitude=3pt, segment length=5.5pt,
                    pre length=3pt, post length=3pt}] (L) -- (R);
  \fill (L) circle (1.4pt);
  \fill (R) circle (1.4pt);
  \draw[-{Stealth[length=5pt]}] (0.3,-0.42) -- (-0.3,-0.42);
  \node[anchor=north, inner sep=2pt] at (0,-0.5) {$\vec q$};
  \node[anchor=north east, inner sep=2pt] at (-2.7,-1.5) {$\vec k,\sigma$};
  \node[anchor=north west, inner sep=2pt] at ( 2.7,-1.5) {$\vec k',\sigma'$};
  \node[anchor=south east, inner sep=2pt] at (-2.7, 1.5) {$\vec k+\vec q,\sigma$};
  \node[anchor=south west, inner sep=2pt] at ( 2.7, 1.5) {$\vec k'-\vec q,\sigma'$};
\end{tikzpicture}
\caption{Vertex corresponding to \cref{eq:h-int}: two electrons exchange a phonon of
momentum $\vec q$.}
\label{fig:vertex}
\end{figure}


Since we know that the momentum transfer is $\vec q$, we know that
$(\vec k,\sigma)\to(\vec k'-\vec q,\sigma')$ is not a valid fermion line.

\subsubsection{Cooper pairs}
\label{subsubsec:cooper-pairs_overview}

In order to examine how this attractive interaction might change the many-body ground
state, one looks at a highly simplified model:
\begin{figure}[ht]
\centering
\begin{tikzpicture}[x=1cm, y=1cm, line width=0.7pt]
  \def\kf{1.9}
  \def\kd{2.45}
  % Debye shell and Fermi sea
  \fill[red!6]   (0,0) circle (\kd);
  \fill[teal!10] (0,0) circle (\kf);
  \draw[red!65!black]  (0,0) circle (\kd);
  \draw[teal!60!black, line width=0.9pt] (0,0) circle (\kf);
  \node[teal!45!black] at (-0.75,0.3) {FS};
  % Fermi momentum
  \fill (0,0) circle (1.4pt);
  \draw[-{Stealth[length=5pt]}] (0,0) -- (-58:\kf);
  \node at (-45:1.15) {$\vec k_F$};
  % width of the shell
  \draw[red!65!black, {Stealth[length=4pt]}-{Stealth[length=4pt]}]
       (192:\kf) -- (192:\kd);
  \node[red!65!black, anchor=east, inner sep=2pt] at (192:3.15) {$\hbar\omega_D$};
  \draw[red!65!black] (192:3.1) -- (192:2.6);
  % two states in the shell, connected by the effective interaction
  \coordinate (A) at (112:2.17);
  \coordinate (B) at (26:2.17);
  \fill (A) circle (1.4pt);
  \fill (B) circle (1.4pt);
  \node[anchor=south east, inner sep=2pt] at (112:2.35) {$\vec k,\sigma$};
  \node[anchor=west, inner sep=2pt]       at (26:2.5)   {$\vec k',\sigma'$};
  \draw[decorate,
        decoration={coil, aspect=0, amplitude=2.5pt, segment length=4.5pt,
                    pre length=2pt, post length=2pt}] (A) to[bend left=38] (B);
  \node[anchor=south, inner sep=2pt] at (69:3.05) {$H_{\mathrm{int}}$};
  \draw (69:3.0) -- (69:2.55);
  \node[anchor=north west, inner sep=2pt] at (1.9,-2.1) {$\vec k$-space};
\end{tikzpicture}
\caption{Only states in a shell of width $\hbar\omega_D$ around the Fermi surface
take part in the effective interaction.}
\label{fig:fermi-surface}
\end{figure}
\begin{itemize}
\item a system of free, non-interacting electrons, i.e.\ the Fermi sea (C\,5.27)
\begin{equation}
|\phi_0\rangle = \prod_{\varepsilon_{\vec k}\le\varepsilon_F,\,\sigma}
  c^\dagger_{\vec k,\sigma}|0\rangle ,
\label{eq:fermi-sea}
\end{equation}
\item the effective interaction is only present for two electrons that are placed right
outside of the Fermi sea with opposite momenta and opposite spins (Cooper pair, singlet
pairing);
\item the interaction is $V=\mathrm{const.}$ if
$|\varepsilon_{\vec k+\vec q}-\varepsilon_{\vec k}|\le\hbar\omega_D$ and zero otherwise.
\end{itemize}
Here ``right outside the Fermi sea'' means that the energies are, cf.\ (C\,5.27),
\begin{equation}
E_F\le\varepsilon_{\vec k},\ \varepsilon_{-\vec k}\le E_F+\hbar\omega_D .
\label{eq:shell}
\end{equation}
The interaction thus reads (C\,5.28)
\begin{equation}
H_{\mathrm{int}} = -\frac V2\sum_{\sigma,\sigma'}\sum_{\vec k,\vec q}
  c^\dagger_{\vec k+\vec q,\sigma}\,c^\dagger_{-\vec k-\vec q,\sigma'}\,
  c_{-\vec k,\sigma'}\,c_{\vec k,\sigma},
  \qquad\text{with }
  |\varepsilon_{\vec k+\vec q}-\varepsilon_{\vec k}|\le\hbar\omega_D .
\label{eq:cooper-int}
\end{equation}
One solves this model by making the variational ansatz (C\,5.26)
\begin{equation}
|\psi_{\sigma,\sigma'}\rangle = \sum_{\vec k}a_{\sigma,\sigma'}(\vec k)\,
  c^\dagger_{\vec k,\sigma}\,c^\dagger_{-\vec k,\sigma'}\,|\phi_0\rangle
\label{eq:cooper-ansatz}
\end{equation}
and minimizing the energy
$E=\langle\psi_{\sigma,\sigma'}|H|\psi_{\sigma,\sigma'}\rangle$ with the constraint that
$\langle\psi_{\sigma,\sigma'}|\psi_{\sigma,\sigma'}\rangle\stackrel{!}{=}1$, which gets us
an upper bound for the energy of the bound state. One finds that (C\,5.35)
\begin{equation}
\Delta E = 2E_F-E_{\mathrm{bs}} \cong -2\hbar\omega_{\vec q}\,e^{-2/(V\rho_0)},
\label{eq:cooper-binding}
\end{equation}
which shows that a bound state (Cooper pair) of the two electrons is
energetically favoured compared to the free, non-interacting state. This holds for all
electrons close to the Fermi edge and therefore implies an instability of the Fermi sea.

\subsubsection{BCS theory}
\label{subsubsec:bcs-theory_overview}

To solve the Hamiltonian and extract the relevant mechanism from it, it is simplified as
follows: Only include interactions between electrons with opposite momenta and opposite
spins, and neglect the $\vec q$-dependence. This yields (C\,5.38)
\begin{equation}
H = \sum_{\vec k,\sigma}\varepsilon_{\vec k}\,c^\dagger_{\vec k\sigma}c_{\vec k\sigma}
  - V\sum_{\vec k,\vec k'}
    c^\dagger_{\vec k'\uparrow}c^\dagger_{-\vec k'\downarrow}
    c_{-\vec k\downarrow}c_{\vec k\uparrow},
\label{eq:bcs-hamiltonian}
\end{equation}
with again $|\varepsilon_{\vec k}-\varepsilon_{\vec k'}|\le\hbar\omega_D$ for
the interaction term. Finally, apply mean-field theory to the pair
creation/destruction operators, (C\,5.41) and (C\,5.42):
\begin{equation}
H_{\mathrm{MF}} = \sum_{\vec k}
  \begin{pmatrix}c^\dagger_{\vec k\uparrow}&c_{-\vec k\downarrow}\end{pmatrix}
  \begin{pmatrix}
    \varepsilon_{\vec k} & -\Delta\\
    -\Delta^* & -\varepsilon_{\vec k}
  \end{pmatrix}
  \begin{pmatrix}c_{\vec k\uparrow}\\c^\dagger_{-\vec k\downarrow}\end{pmatrix}
  \;\pm\;\mathrm{const.},
\label{eq:h-mf}
\end{equation}
\begin{equation}
\Delta = V\sum_{\vec k}\langle c_{-\vec k\downarrow}c_{\vec k\uparrow}\rangle
\qquad(\text{possibly }\varepsilon_{\vec k}\rightsquigarrow\varepsilon_{\vec k}-\mu).
\label{eq:gap-def}
\end{equation}

\paragraph{Bogoliubov transformation.}
Introduce a unitary transformation that diagonalizes
\begin{equation*}
B = \left(\begin{matrix}\varepsilon_{\vec k}&-\Delta\\
-\Delta^*&-\varepsilon_{\vec k}\end{matrix}\right)     
\end{equation*}
and thus leads to an
$H_{\mathrm{MF}}$ that only includes terms of the form
$\alpha^\dagger_{\vec k}\alpha_{\vec k}$ and $\beta^\dagger_{\vec k}\beta_{\vec k}$:
\begin{equation}
H_{\mathrm{MF}} = \sum_{\vec k}
  \underbrace{
  \begin{pmatrix}c^\dagger_{\vec k\uparrow}&c_{-\vec k\downarrow}\end{pmatrix}
  U_{\vec k}}_{\text{new operators}}\;
  \underbrace{U^\dagger_{\vec k}
  \begin{pmatrix}
    \varepsilon_{\vec k} & -\Delta\\ -\Delta^* & -\varepsilon_{\vec k}
  \end{pmatrix}
  U_{\vec k}}_{\text{diagonal matrix}}\;
  \underbrace{U^\dagger_{\vec k}
  \begin{pmatrix}c_{\vec k\uparrow}\\c^\dagger_{-\vec k\downarrow}\end{pmatrix}
  }_{(\text{new operators})^\dagger}.
\label{eq:bogoliubov}
\end{equation}
These new operators will be fermionic, since $U$ is unitary, so that
$\langle c_{-\vec k\downarrow}c_{\vec k\uparrow}\rangle\rightsquigarrow
f(\alpha_{\vec k},\beta_{\vec k},\alpha^\dagger_{\vec k},\beta^\dagger_{\vec k})$ and one
can then use $\langle\alpha^\dagger_{\vec k}\alpha_{\vec k}\rangle=f(E_F)$ and
$\langle\beta_{\vec k}\alpha_{\vec k}\rangle=0$ etc. The elements of $U$ can be determined
by diagonalizing $B$ and requiring that the eigenvectors are normalized to one.

\paragraph{Result.}
This yields
\begin{align}
H_{\mathrm{eff}} &= \sum_{\vec k}E_{\vec k}
  \big(\alpha^\dagger_{\vec k}\alpha_{\vec k}
     + \beta^\dagger_{\vec k}\beta_{\vec k}\big) + \mathrm{const.},
  &&\text{(C\,5.46)}
\label{eq:h-eff}\\
E_{\vec k} &= \big(\varepsilon_{\vec k}^2+|\Delta|^2\big)^{1/2}\ \ge\ 0 ,
  &&\text{(C\,5.47)}
\label{eq:e-k}
\end{align}
and for $\Delta$ one arrives at the BCS self-consistency equation (C\,5.50)
\begin{equation}
\Delta = \frac{V\Delta}{2}\sum_{\vec k}\frac{1}{E_{\vec k}}
  \tanh\!\left(\frac{\beta E_{\vec k}}{2}\right).
\label{eq:gap-selfconsistency}
\end{equation}
Obviously $\Delta=0$ is a solution. For $\Delta\neq0$ one obtains (C\,5.51)
\begin{equation}
1 = \frac V2\sum_{\vec k, |\varepsilon_{\vec k}|<\hbar\omega_D}
  \big(\varepsilon_{\vec k}^2+|\Delta|^2\big)^{-1/2}
  \tanh\!\left(
    \frac{\beta\big(\varepsilon_{\vec k}^2+|\Delta|^2\big)^{1/2}}{2}\right).
\label{eq:gap-equation}
\end{equation}

\paragraph{$\Delta$ as an order parameter.}
Actually, $\Delta$ is an order parameter for the superconducting phase, because of the
following:
\begin{itemize}
\item For a normal metal
$\langle c_{\vec k\uparrow}c_{-\vec k\downarrow}\rangle=0\ \forall\,\vec k$ and thus
$\Delta = V\sum_{\vec k}\langle c_{-\vec k\downarrow}c_{\vec k\uparrow}\rangle = 0$.
\item But if $\Delta\neq0$, it means that we are not in a metallic phase and that there
are Cooper pairs present.
\end{itemize}

\subsubsection{The BCS ground state}
\label{subsubsec:bcs-ground-state_overview}

We have $H=\sum_{\vec k}E_{\vec k}(\alpha^\dagger_{\vec k}\alpha_{\vec k}
+\beta^\dagger_{\vec k}\beta_{\vec k})$ with $E_{\vec k}\ge0$, cf.\ \cref{eq:h-eff}, where,
following (C\,5.43),
\begin{equation}
\alpha_{\vec k} = u_{\vec k}c_{\vec k\uparrow}-v_{\vec k}c^\dagger_{-\vec k\downarrow}
\qquad\text{and}\qquad
\beta_{\vec k} = u_{\vec k}c_{-\vec k\downarrow}+v_{\vec k}c^\dagger_{\vec k\uparrow}.
\label{eq:bogoliubov-ops}
\end{equation}
Since $\alpha^{(\dagger)}_{\vec k}$ and $\beta^{(\dagger)}_{\vec k}$ are
fermionic, the eigenvalues of $\alpha^\dagger_{\vec k}\alpha_{\vec k}$ (and
$\beta^\dagger_{\vec k}\beta_{\vec k}$) are $0,1$. Hence the ground state of $H$ does not
contain any $\alpha$ or $\beta$ quasiparticles, i.e.
\begin{equation}
\alpha_{\vec k}|\mathrm{BCS}\rangle = \beta_{\vec k}|\mathrm{BCS}\rangle = 0 .
\label{eq:quasiparticle-vacuum}
\end{equation}
This is satisfied by
\begin{equation}
|\mathrm{BCS}\rangle
 = \prod_{\vec k}\big(u_{\vec k}
   + v_{\vec k}c^\dagger_{\vec k\uparrow}c^\dagger_{-\vec k\downarrow}\big)|0\rangle
 = \bigotimes_{\vec k}\big(u_{\vec k}
   + v_{\vec k}c^\dagger_{\vec k\uparrow}c^\dagger_{-\vec k\downarrow}\big)
   |0\rangle_{\vec k}.
\label{eq:bcs-state}
\end{equation}


\subsection{Details} \label{subsec:froehlich-to-bcs_details}
\subsubsection{Preface}

These are supplementary notes to \emph{Solid State Theory, Volume 2} by Gerd Czycholl.
Since the derivations there are quite lengthy, only those steps are presented here that
clarify what is done in the corresponding chapters. Numbers in parentheses such as
(5.12) refer to equations of that book.

\subsubsection{Attractive electron-electron interaction through the electron-phonon
mechanism (Czycholl 5.2)}
\label{subsubsec:bcs-el-ph_details}

\paragraph{Two small observations.}
\begin{itemize}
\item (5.12): use that $S$ is linear in $H_1$ to discard $[[H_1,S],S]$.
\item (5.15): use that $\langle m|H_1|m\rangle=0$, since $H_1$ changes the number of
      phonons and thus $H_1|m\rangle\perp|m\rangle$.
\end{itemize}

\paragraph{The transformed Hamiltonian.}
Starting from (5.13) and inserting resolutions of the identity,
\begin{align}
H_T &= H_0+\frac{i}{2}\big(H_1S-SH_1\big)\\
&= H_0+\frac{i}{2}\sum_{\tilde n,n,m}\Big(
   |\tilde n\rangle\langle\tilde n|H_1|m\rangle\langle m|S|n\rangle\langle n|
 - |\tilde n\rangle\langle\tilde n|S|m\rangle\langle m|H_1|n\rangle\langle n|\Big)\\
&\overset{(5.15)}{=} H_0-\frac{1}{2}\sum_{\tilde n,n,m}\Big(
   |\tilde n\rangle\langle\tilde n|H_1|m\rangle\langle m|H_1|n\rangle\langle n|
   \frac{1}{E_m-E_n}\nonumber\\
&\hspace{7em}
 - |\tilde n\rangle\langle\tilde n|H_1|m\rangle\langle m|H_1|n\rangle\langle n|
   \frac{1}{E_{\tilde n}-E_m}\Big),
\end{align}
where the matrix elements $\langle\tilde n|H_1|m\rangle$ and $\langle m|H_1|n\rangle$ are
scalars and can be pulled out of the operator structure, which yields (5.16).

\paragraph{Restriction to a fixed phonon configuration.}
Since we want to derive a Hamiltonian for the electron-electron interaction only, we
restrict $|\tilde n\rangle$ and $|n\rangle$ in (5.16) to states with the same phonon
configuration, such that $H_T$ does not mix states with different phonon configurations.

Let $|n\rangle$ be some many-particle state and assume that it contains one electron in
the state $|k,\sigma\rangle$ and one in $|k',\sigma'\rangle$. The two elementary
processes are
\begin{align}
H_{1,1}|k,\sigma\rangle &= \sqrt{n_q}\,|k+q,\sigma,n_q-1\rangle,
  &&\text{by destroying a phonon,} \tag{$\ast_1$}\\
H_{1,2}|k,\sigma\rangle &= \sqrt{n_q+1}\,|k-q,\sigma,n_q+1\rangle,
  &&\text{by creating a phonon.} \tag{$\ast_2$}
\end{align}
Since $|\tilde n\rangle$ and $|n\rangle$ are restricted to states with the same phonon
number, the created (destroyed) phonon with momentum $q$ has to be destroyed (created)
again:
\begin{align}
H_{1,1}\sqrt{n_q}\,|k+q,\sigma,n_q-1\rangle
  &= n_q\,|k+q,\sigma;\,k'-q,\sigma';\,n_q\rangle,\\
H_{1,1}\sqrt{n_q+1}\,|k-q,\sigma,n_q+1\rangle
  &= (n_q+1)\,|k-q,\sigma;\,k'+q,\sigma';\,n_q\rangle.
\end{align}
All other terms vanish, since their phonon configurations differ.

\paragraph{Energy denominators.}
Therefore we have
\begin{align}
(\ast_1):&\quad E_m-E_n=\varepsilon_{k+q}-\varepsilon_k-\hbar\omega_q,
  &&E_{\tilde n}-E_m=\varepsilon_{k'-q}-\varepsilon_{k'}+\hbar\omega_q,
  &&(5.19)\\
(\ast_2):&\quad E_m-E_n=\varepsilon_{k-q}-\varepsilon_k+\hbar\omega_q,
  &&E_{\tilde n}-E_m=\varepsilon_{k'+q}-\varepsilon_{k'}-\hbar\omega_q.
  &&(5.20)
\end{align}
Inserting this into (5.16) gives
\begin{equation}
\begin{split}
H_T=-\frac{1}{2}\sum_{n,\tilde n}|\tilde n\rangle\langle n|\sum_m
\sum_{\substack{k,k',q\\ \sigma,\sigma'}}\Big(&
\langle\tilde n|M(-q)c^\dagger_{k'-q,\sigma'}c_{k',\sigma'}b^\dagger_q|m\rangle
\langle m|M(q)c^\dagger_{k+q,\sigma}c_{k,\sigma}b_q|n\rangle\\
&\quad\times\Big(\frac{1}{\varepsilon_{k+q}-\hbar\omega_q-\varepsilon_k}
-\frac{1}{\varepsilon_{k'-q}-\varepsilon_{k'}+\hbar\omega_q}\Big)\\
+\;&\langle\tilde n|M(q)c^\dagger_{k'+q,\sigma'}c_{k',\sigma'}b_q|m\rangle
\langle m|M(-q)c^\dagger_{k-q,\sigma}c_{k,\sigma}b^\dagger_q|n\rangle\\
&\quad\times\Big(\frac{1}{\varepsilon_{k-q}-\varepsilon_k+\hbar\omega_q}
-\frac{1}{\varepsilon_{k'+q}-\varepsilon_{k'}-\hbar\omega_q}\Big)\Big),
\end{split}
\end{equation}
where the first summand is the $(\ast_1)$-process and the second one the
$(\ast_2)$-process.

Using $M^*(q)=M(-q)$ and that all the above terms are independent of $n$, $m$, and
$\tilde n$ yields (5.21). Using furthermore that
\begin{equation}
b^\dagger_q b_q|\dots n_q\dots\rangle=n_q|\dots n_q\dots\rangle,\qquad
b_q b^\dagger_q|\dots n_q\dots\rangle=(n_q+1)|\dots n_q\dots\rangle
\end{equation}
results in (5.21), i.e.\ $\hat H_T(b_q,b^\dagger_q,\dots)\rightsquigarrow
\hat H_T(n_q,\dots)$ in the phonon number basis.

\paragraph{Combining the two processes.}
We have
\begin{equation}
\frac{1}{\varepsilon_{k+q}-\hbar\omega_q-\varepsilon_k}
=\frac{(\varepsilon_{k+q}-\varepsilon_k)+\hbar\omega_q}
      {(\varepsilon_{k+q}-\varepsilon_k)^2-\hbar^2\omega_q^2}
\end{equation}
and, relabelling $k\leftrightarrow k'$ and $q\to-q$,
\begin{equation}
\sum_{k,k',q}\frac{1}{\varepsilon_{k'-q}-\varepsilon_{k'}+\hbar\omega_q}\,
c^\dagger_{k+q}c_k c^\dagger_{k'-q}c_{k'}
=\sum_{k,k',q}\frac{(\varepsilon_{k+q}-\varepsilon_k)-\hbar\omega_q}
      {(\varepsilon_{k+q}-\varepsilon_k)^2-\hbar^2\omega_q^2}\,
c^\dagger_{k'-q}c_{k'}c^\dagger_{k+q}c_k ,
\end{equation}
where
\begin{align}
\sum_{k,k',q}c^\dagger_{k'-q}c_{k'}c^\dagger_{k+q}c_k
&=\sum_{k,k',q}c^\dagger_{k'-q}\big(\delta_{k',k+q}-c^\dagger_{k+q}c_{k'}\big)c_k\\
&=\sum_{k,q}c^\dagger_k c_k-\sum_{k,k',q}c^\dagger_{k+q}c^\dagger_{k'-q}c_k c_{k'}\\
&=\sum_{k,q}n_k-\sum_{k,k',q}c^\dagger_{k+q}
   \big(\delta_{k,k'-q}-c_k c^\dagger_{k'-q}\big)c_{k'}\\
&=\underbrace{\sum_{k,q}n_k-\sum_{k,q}c^\dagger_{k'}c_{k'}}_{=\,0}
  +\sum_{k,k',q}c^\dagger_{k+q}c_k c^\dagger_{k'-q}c_{k'}\\
&=\sum_{k,k',q}c^\dagger_{k+q}c_k c^\dagger_{k'-q}c_{k'} .
\end{align}

We know (Bruus \& Flensberg, Sec.~3.8) that $\omega_{-q}=\omega_q$ and $M(-q)=M(q)$.
Furthermore, since we are only interested in the electronic process, we can assume that
there is no net phononic motion within the system and thus $n_q=n_{-q}$.
Hence the first term of (5.21) equals the first term of (5.22), analogously for the
second term, and similarly (5.22) $\rightsquigarrow$ (5.23).

\subsubsection{Cooper pairs (Czycholl 5.3)}
\label{subsubsec:bcs-cooper-pairs_details}

\paragraph{Setup.}
Simplify: look at the Fermi sea plus two electrons just above the Fermi energy, and only
let them interact via $H_T$. The variational ansatz (5.26), for vanishing total pair
momentum $K=k+k'\overset{!}{=}0$, is
\begin{equation}
|\psi_{\sigma,\sigma'}\rangle=\sum_k a_{\sigma,\sigma'}(k)\,
c^\dagger_{k,\sigma}c^\dagger_{-k,\sigma'}|\phi_0\rangle ,
\end{equation}
with the sum over $k$ such that $E_F\le\varepsilon_k,\varepsilon_{k'}\le
E_F+\hbar\omega_D$, cf.\ (5.27). The simplified interaction acts only outside the Fermi
sea ($\ast$), otherwise $V=0$:
\begin{equation}
H_T=-\frac{V}{2}\sum_{\sigma,\sigma'}\sum_{k,q}
c^\dagger_{k+q,\sigma}c^\dagger_{-k-q,\sigma'}c_{-k,\sigma'}c_{k,\sigma}.
\qquad (5.28)
\end{equation}

\paragraph{Action on the ansatz.}
\begin{align}
H_T|\psi_{\sigma_1\sigma_2}\rangle
&=-\frac{V}{2}\sum_{\sigma,\sigma'}\sum_{k,q,k'}
c^\dagger_{k+q,\sigma}c^\dagger_{-k-q,\sigma'}c_{-k,\sigma'}c_{k,\sigma}\,
a_{\sigma_1\sigma_2}(k')\,c^\dagger_{k',\sigma_1}c^\dagger_{-k',\sigma_2}|\phi_0\rangle\\
&=-\frac{V}{2}\sum_{\sigma,\sigma'}\sum_{k,q,k'}a_{\sigma_1\sigma_2}(k')\,
c^\dagger_{k+q,\sigma}c^\dagger_{-k-q,\sigma'}c_{-k,\sigma'}c_{k,\sigma}
\,|k',\sigma_1;-k',\sigma_2\rangle\\
&=-\frac{V}{2}\sum_{\sigma,\sigma'}\sum_{k,q,k'}a_{\sigma_1\sigma_2}(k')\Big(
 \delta_{k,k'}\delta_{\sigma,\sigma_1}\delta_{\sigma',\sigma_2}
 |k+q,\sigma_1;-k-q,\sigma_2\rangle\nonumber\\
&\hspace{10em}
 +\delta_{k,-k'}\delta_{\sigma,\sigma_2}\delta_{\sigma',\sigma_1}
 |-k-q,\sigma_1;k+q,\sigma_2\rangle\Big)\\
&=-\frac{V}{2}\sum_{k,q}\Big(a_{\sigma_1\sigma_2}(k)\,
 |k+q,\sigma_1;-k-q,\sigma_2\rangle
 +a_{\sigma_1\sigma_2}(-k)\,|-k-q,\sigma_1;k+q,\sigma_2\rangle\Big),
\end{align}
where the projection onto $|k',\sigma_1;-k',\sigma_2\rangle$ is allowed because, by
($\ast$), only states outside the Fermi sea matter for the interaction. Taking the
expectation value,
\begin{align}
\langle\psi_{\sigma_1\sigma_2}|H_T|\psi_{\sigma_1\sigma_2}\rangle
&=-\frac{V}{2}\sum_{k,q,k'}a^*_{\sigma_1\sigma_2}(k')\Big\{
 a_{\sigma_1\sigma_2}(k)\,
 \langle k',\sigma_1;-k',\sigma_2|k+q,\sigma_1;-k-q,\sigma_2\rangle\nonumber\\
&\hspace{8em}
 +a_{\sigma_1\sigma_2}(-k)\,
 \langle k',\sigma_1;-k',\sigma_2|-k-q,\sigma_1;k+q,\sigma_2\rangle\Big\}\\
&=-\frac{V}{2}\sum_{k,q,k'}a^*_{\sigma_1\sigma_2}(k')\Big\{
 a_{\sigma_1\sigma_2}(k)\,\delta_{k',k+q}
 +a_{\sigma_1\sigma_2}(-k)\,\delta_{k',-k-q}\Big\}\\
&=-\frac{V}{2}\sum_{k,q}\Big(a^*_{\sigma_1\sigma_2}(k+q)a_{\sigma_1\sigma_2}(k)
 +a^*_{\sigma_1\sigma_2}(-k-q)a_{\sigma_1\sigma_2}(-k)\Big)\\
&=-V\sum_{k,q}a^*_{\sigma_1\sigma_2}(k+q)a_{\sigma_1\sigma_2}(k),
\end{align}
where in the last step $k\to-k$, $q\to-q$ was substituted in the second term.

\paragraph{The kinetic part.}
Furthermore, by (5.17),
\begin{equation}
H_0=\sum_{k,\sigma}\varepsilon_k n_{k,\sigma}+\sum_q\hbar\omega_q b^\dagger_q b_q .
\end{equation}
Since the ansatz $|\psi_{\sigma\sigma'}\rangle$ is purely electronic,
$\sum_q\hbar\omega_q b^\dagger_q b_q$ will not contribute. Splitting the electronic sum
into states inside and outside the Fermi sea,
\begin{align}
\langle\psi_{\sigma,\sigma'}|H_0|\psi_{\sigma,\sigma'}\rangle
&=\langle\psi_{\sigma\sigma'}|\Big(\sum_{k,\tilde\sigma}
  \varepsilon_k n_{k,\tilde\sigma}\Big)
  \Big(\sum_{k'}a_{\sigma,\sigma'}(k')c^\dagger_{k',\sigma}
  c^\dagger_{-k',\sigma'}|\phi_0\rangle\Big)\\
&=\langle\psi_{\sigma\sigma'}|\Big(\sum_{\substack{k\le k_F\\ \tilde\sigma}}
  \varepsilon_k\hat n_{k,\tilde\sigma}
  +\sum_{\substack{k>k_F\\ \tilde\sigma}}\varepsilon_k\hat n_{k,\tilde\sigma}\Big)
  \Big(\sum_{k'}a_{\sigma,\sigma'}(k')c^\dagger_{k',\sigma}
  c^\dagger_{-k',\sigma'}|\phi_0\rangle\Big)\\
&=\langle\psi_{\sigma\sigma'}|\sum_{\substack{k\le k_F\\ \tilde\sigma}}
  \varepsilon_k n_{k,\tilde\sigma}|\psi_{\sigma\sigma'}\rangle
  +\sum_{k,k'}a_{\sigma,\sigma'}(k')\big(\varepsilon_{k'}\delta_{k,k'}
  \delta_{\sigma,\tilde\sigma}\nonumber\\
&\hspace{10em}
  +\varepsilon_{k'}\delta_{k,-k'}\delta_{\sigma',\tilde\sigma}\big)
  c^\dagger_{k',\sigma}c^\dagger_{-k',\sigma'}|\phi_0\rangle\\
&=\Big(\langle\phi_0|\sum_k c_{-k,\sigma'}c_{k,\sigma}a^*_{\sigma,\sigma'}(k)\Big)
  \Big(\sum_{\substack{k\le k_F\\ \tilde\sigma}}\varepsilon_k n_{k,\tilde\sigma}
  |\psi_{\sigma\sigma'}\rangle\nonumber\\
&\hspace{8em}
  +\sum_{k'}a_{\sigma,\sigma'}(k')(\varepsilon_{k'}+\varepsilon_{-k'})
  c^\dagger_{k',\sigma}c^\dagger_{-k',\sigma'}|\phi_0\rangle\Big)\\
&=\sum_k|a_{\sigma,\sigma'}(k)|^2\Big((\varepsilon_k+\varepsilon_{-k})
  +\sum_{\substack{k\le k_F\\ \tilde\sigma}}\varepsilon_k n_{k,\tilde\sigma}\Big)\\
&=\sum_k|a_{\sigma,\sigma'}(k)|^2(\varepsilon_k+\varepsilon_{-k})
  +c\sum_k|a_{\sigma\sigma'}(k)|^2 ,
\end{align}
where $c$ collects the (constant) energy of the filled Fermi sea. From
$\langle\psi_{\sigma\sigma'}|\psi_{\sigma\sigma'}\rangle\overset{!}{=}1$ it follows that
$\sum_k|a_{\sigma\sigma'}(k)|^2\overset{!}{=}1$, cf.\ (5.29), and therefore
\begin{equation}
H_0=\sum_k|a_{\sigma,\sigma'}(k)|^2(\varepsilon_k+\varepsilon_{-k}),
\qquad\Longrightarrow\qquad H=H_0+H_T=(5.29).
\end{equation}
The step $\varepsilon_k=\varepsilon_{-k}$ needed for (5.30) is true due to Kramers'
theorem; see e.g.\ Sec.~5.3 of Czycholl~I.

\paragraph{When does the gap equation degenerate?}
For (5.32) it was assumed that $C\equiv\sum_k a_{\sigma\sigma'}(k)\neq0$. Let us examine
when $C$ is actually zero. We have
\begin{equation}
a_{\sigma\sigma'}(k)=\frac{V\sum_{k'}a_{\sigma\sigma'}(k')}{2\varepsilon_k-\lambda},
\qquad (5.31)
\end{equation}
where $\varepsilon_k=\varepsilon_{-k}$ by Kramers' theorem, and hence
\begin{equation}
a_{\sigma\sigma'}(-k)=a_{\sigma\sigma'}(k). \tag{$\ast$}
\end{equation}
Moreover, for equal spins,
\begin{align}
|\psi_{\sigma\sigma}\rangle
&=\sum_k a_{\sigma,\sigma}(k)c^\dagger_{k,\sigma}c^\dagger_{-k,\sigma}|\phi_0\rangle\\
&=\sum_k a_{\sigma,\sigma}(k)\big(\delta_{k,-k}
  -c^\dagger_{-k,\sigma}c^\dagger_{k,\sigma}\big)|\phi_0\rangle\\
&=-\sum_k a_{\sigma,\sigma}(k)c^\dagger_{-k,\sigma}c^\dagger_{k,\sigma}|\phi_0\rangle\\
&=-\sum_k a_{\sigma,\sigma}(-k)c^\dagger_{k,\sigma}c^\dagger_{-k,\sigma}|\phi_0\rangle\\
&=-\sum_k a_{\sigma,\sigma}(k)c^\dagger_{k,\sigma}c^\dagger_{-k,\sigma}|\phi_0\rangle
 =-|\psi_{\sigma\sigma}\rangle .
\end{align}
Here $\delta_{k,-k}$ drops out because $\sum_k$ does not include momenta inside the Fermi
sea and thus $k\neq0$; in the fourth line $k\to-k$ was substituted, and in the fifth line
($\ast$) was used. Consequently
\begin{equation}
\langle\psi_{\sigma\sigma}|\psi_{\sigma\sigma}\rangle
=-\langle\psi_{\sigma\sigma}|\psi_{\sigma\sigma}\rangle
\;\Longrightarrow\;\langle\psi_{\sigma\sigma}|\psi_{\sigma\sigma}\rangle=0
\;\Longrightarrow\;\sum_k|a_{\sigma\sigma}(k)|^2=0
\;\Longrightarrow\;a_{\sigma\sigma}(k)=0,
\end{equation}
so $\sum_k a_{\sigma\sigma}(k)=0$. Thus $\sigma=\sigma'$ yields $C=0$, cf.\ (5.35), and
Cooper pairs consist of electrons with opposite spins.

Note that if $\psi(r_1-r_2)$ is antisymmetric, then $\psi(0)=0$ and $C=0$ as well, so the
derivation of (5.35) is different in that case.

\paragraph{Position representation.}
Project the state
$|\psi_{\sigma,\sigma'}\rangle=\sum_k a_{\sigma,\sigma'}(k)c^\dagger_{k,\sigma}
c^\dagger_{-k,\sigma'}|\phi_0\rangle$ onto the position basis. Since $H$ acts trivially
on the spin space ($\sigma\to\sigma$, $\sigma'\to\sigma'$), one can factorize the wave
function:
\begin{equation}
\langle \vec r_1,\vec r_2|\psi_{\sigma,\sigma'}\rangle
=\psi_{\sigma,\sigma'}(\vec r_1,\vec r_2)
=\psi(\vec r_1,\vec r_2)\,\chi(\sigma,\sigma'),
\qquad (5.37)
\end{equation}
with
\begin{equation}
\psi(\vec r_1,\vec r_2)=\frac{1}{V}\sum_k a(k)\,
e^{i\vec k\cdot\vec r_1}e^{-i\vec k\cdot\vec r_2}=\psi(\vec r_1-\vec r_2),
\qquad (5.36)
\end{equation}
using $\langle\vec r|\vec k\rangle=\frac{1}{\sqrt V}e^{i\vec k\cdot\vec r}$.

\subsubsection{BCS theory (Czycholl 5.4)}
\label{subsubsec:bcs-theory_details}

\paragraph{Mean-field recap.}
Write $\hat n_\alpha=\langle n_\alpha\rangle+\delta\hat n_\alpha$ and assume that
$\delta\hat n_\alpha$ is small, such that second-order fluctuations can be dropped:
\begin{align}
\hat n_\alpha\hat n_\beta
&=\big(\langle\hat n_\alpha\rangle+\delta\hat n_\alpha\big)
  \big(\langle\hat n_\beta\rangle+\delta\hat n_\beta\big)\\
&\overset{\delta\hat n_\alpha\delta\hat n_\beta\approx0}{=}
  \langle\hat n_\alpha\rangle\langle\hat n_\beta\rangle
  +\langle\hat n_\alpha\rangle\big(\hat n_\beta-\langle\hat n_\beta\rangle\big)
  +\langle\hat n_\beta\rangle\big(\hat n_\alpha-\langle\hat n_\alpha\rangle\big)\\
&=\langle\hat n_\alpha\rangle\hat n_\beta+\langle\hat n_\beta\rangle\hat n_\alpha
  -\langle\hat n_\alpha\rangle\langle\hat n_\beta\rangle .
\end{align}
Now apply this to $c^{(\dagger)}_{-k\downarrow}c^{(\dagger)}_{k\uparrow}$ instead of
$\hat n_\alpha$ to derive a mean-field version of the BCS Hamiltonian, which gives
(5.40).

\paragraph{Nambu form of the effective Hamiltonian.}
Rewrite $H_{\mathrm{eff}}$ as (following Becca \& Sorella, and the Superconductivity
IMPRS UFAST 2017 course of Michael Sentef)
\begin{align}
\tilde H_{\mathrm{eff}}
&=\sum_k\begin{pmatrix}c^\dagger_{k\uparrow}&c_{-k\downarrow}\end{pmatrix}
  \begin{pmatrix}\varepsilon_k&-\Delta\\-\Delta^*&-\varepsilon_k\end{pmatrix}
  \begin{pmatrix}c_{k\uparrow}\\c^\dagger_{-k\downarrow}\end{pmatrix}\\
&=\sum_k\begin{pmatrix}c^\dagger_{k\uparrow}&c_{-k\downarrow}\end{pmatrix}
  \begin{pmatrix}\varepsilon_k c_{k\uparrow}-\Delta c^\dagger_{-k\downarrow}\\
  -\Delta^* c_{k\uparrow}-\varepsilon_k c^\dagger_{-k\downarrow}\end{pmatrix}\\
&=\sum_k\varepsilon_k c^\dagger_{k\uparrow}c_{k\uparrow}
  -\Delta c^\dagger_{k\uparrow}c^\dagger_{-k\downarrow}
  -\Delta^* c_{-k\downarrow}c_{k\uparrow}
  -\varepsilon_k \underbrace{c_{-k\downarrow}c^\dagger_{-k\downarrow}}
   _{=\,1-c^\dagger_{-k\downarrow}c_{-k\downarrow}}\\
&=\sum_{k,\sigma}\varepsilon_k c^\dagger_{k\sigma}c_{k\sigma}
  -\Delta^*\sum_k c_{-k\downarrow}c_{k\uparrow}
  -\Delta\sum_k c^\dagger_{k\uparrow}c^\dagger_{-k\downarrow}
  -\sum_k\varepsilon_k\\
&=H_{\mathrm{eff}}-\underbrace{\frac{|\Delta|^2}{V}-\sum_k\varepsilon_k}
   _{\text{const.}},\qquad (5.41)
\end{align}
Set $M\equiv\begin{pmatrix}\varepsilon_k&-\Delta\\-\Delta^*&-\varepsilon_k\end{pmatrix}$
and diagonalize it, such that $H_{\mathrm{eff}}$ becomes diagonal --- this is a unitary
transformation.

\paragraph{Diagonalizing the $2\times2$ matrix $M$.}
Eigenvalues:
\begin{align}
\det\begin{pmatrix}\varepsilon_k-\lambda&-\Delta\\
 -\Delta^*&-\varepsilon_k-\lambda\end{pmatrix}
&=-(\varepsilon_k-\lambda)(\varepsilon_k+\lambda)-|\Delta|^2\\
&=-\big(\lambda^2-\varepsilon_k^2-|\Delta|^2\big)\overset{!}{=}0\\
\Longrightarrow\quad \lambda_{1,2}
&=\pm\sqrt{\varepsilon_k^2+|\Delta|^2}\equiv E_k^\pm\in\mathbb{R}. \tag{$\ast_1$}
\end{align}
Eigenvectors: from
\begin{equation}
\begin{pmatrix}\varepsilon_k-E_k^\pm&-\Delta\\
 -\Delta^*&-\varepsilon_k-E_k^\pm\end{pmatrix}
\begin{pmatrix}x_\pm\\y_\pm\end{pmatrix}\overset{!}{=}0
\end{equation}
the first row (I) gives
\begin{equation}
(\varepsilon_k-E_k^\pm)x_\pm-\Delta y_\pm=0
\quad\Longleftrightarrow\quad
y_\pm=\frac{\varepsilon_k-E_k^\pm}{\Delta}\,x_\pm, \tag{I$'$}
\end{equation}
while the second row (II) gives, using (I$'$) and
$(E_k^\pm)^2=\varepsilon_k^2+|\Delta|^2$,
\begin{equation}
\Delta^*x_\pm+(\varepsilon_k+E_k^\pm)y_\pm
\overset{\text{(I}'\text{)}}{=}
\Big(\Delta^*+\frac{\varepsilon_k^2-(E_k^\pm)^2}{\Delta}\Big)x_\pm=0 ,
\end{equation}
i.e.\ it is satisfied identically. Hence $x_\pm$ can be chosen arbitrarily. \hfill($\odot$)

\paragraph{Normalization.}
Requiring $|x_\pm|^2+|y_\pm|^2\overset{!}{=}1$ and using (I$'$),
\begin{equation}
|x_\pm|^2\Big(1+\frac{(\varepsilon_k-E_k^\pm)^2}{|\Delta|^2}\Big)
=|x_\pm|^2\,\frac{|\Delta|^2+(\varepsilon_k-E_k^\pm)^2}{|\Delta|^2}\overset{!}{=}1,
\end{equation}
where
\begin{align}
(\varepsilon_k-E_k^\pm)^2
&=\varepsilon_k^2-2\varepsilon_k E_k^\pm+(E_k^\pm)^2\\
&=\varepsilon_k^2-2\varepsilon_k E_k^\pm+\varepsilon_k^2+|\Delta|^2\\
&=2\varepsilon_k^2-2\varepsilon_k E_k^\pm+|\Delta|^2 .
\end{align}
Therefore
\begin{align}
|x_\pm|^2
&=\frac{|\Delta|^2}{2\big(\varepsilon_k^2-\varepsilon_k E_k^\pm+|\Delta|^2\big)}
 \overset{(\ast_1)}{=}\frac{1}{2}\,
 \frac{|\Delta|^2}{(E_k^\pm)^2-\varepsilon_k E_k^\pm}
 \overset{(\ast_1)}{=}\frac{1}{2}\,
 \frac{(E_k^\pm)^2-\varepsilon_k^2}{E_k^\pm\big(E_k^\pm-\varepsilon_k\big)}\\
&=\frac{1}{2}\,\frac{E_k^\pm+\varepsilon_k}{E_k^\pm}
 =\frac{1}{2}\Big(1+\frac{\varepsilon_k}{E_k^\pm}\Big)
 =\frac{1}{2}\Big(1\pm\frac{\varepsilon_k}{E_k}\Big),\\
|y_\pm|^2&=1-|x_\pm|^2
 =\frac{1}{2}\Big(1-\frac{\varepsilon_k}{E_k^\pm}\Big)
 =\frac{1}{2}\Big(1\mp\frac{\varepsilon_k}{E_k}\Big).
\tag{A}
\end{align}

\paragraph{Fixing the free phase.}
Choose $x_\pm$ to be some real value, which is allowed by ($\odot$). Then (A) gives
$x_\pm^2=y_\mp^2$, i.e.\ $x_\pm=\pm y_\mp$, which is still too many solutions. Set
\begin{equation}
x_+=y_-\equiv u_k,\qquad x_-=-y_+\equiv v_k .
\end{equation}
Since $M^\dagger_k=M_k$, we know that $M_k$ can be diagonalized by a unitary $U_k$:
\begin{equation}
M=\begin{pmatrix}\varepsilon_k&-\Delta\\-\Delta^*&-\varepsilon_k\end{pmatrix}
=U^\dagger_k\begin{pmatrix}E_k^+&0\\0&E_k^-\end{pmatrix}U_k,
\qquad
U_k=\begin{pmatrix}x^+&x^-\\y^+&y^-\end{pmatrix}
=\begin{pmatrix}u_k&v_k\\-v_k&u_k\end{pmatrix},
\end{equation}
which is obviously unitary.

\paragraph{The Bogoliubov transformation.}
Inserting this into the Hamiltonian yields
\begin{align}
\tilde H_{\mathrm{eff}}
&=\sum_k\begin{pmatrix}c^\dagger_{k\uparrow}&c_{-k\downarrow}\end{pmatrix}
  U_kU^\dagger_k
  \begin{pmatrix}\varepsilon_k&-\Delta\\-\Delta^*&-\varepsilon_k\end{pmatrix}
  U_kU^\dagger_k
  \begin{pmatrix}c_{k\uparrow}\\c^\dagger_{-k\downarrow}\end{pmatrix}
  -\sum_k\varepsilon_k\\
&=\sum_k\begin{pmatrix}\alpha^\dagger_k&\beta_k\end{pmatrix}
  \begin{pmatrix}E_k^+&0\\0&E_k^-\end{pmatrix}
  \begin{pmatrix}\alpha_k\\\beta^\dagger_k\end{pmatrix}-\sum_k\varepsilon_k\\
&=\sum_k E_k^+\alpha^\dagger_k\alpha_k+E_k^-\beta_k\beta^\dagger_k
  -\sum_k\varepsilon_k\\
&=\sum_k E_k^+\big(\alpha^\dagger_k\alpha_k+\beta^\dagger_k\beta_k\big)
  -\underbrace{\sum_k\big(\varepsilon_k+E_k^+\big)}_{\text{const.}},
  \qquad\sim(5.46)
\end{align}
where $\{\beta_k,\beta^\dagger_k\}=1$ and $E_k^-=-E_k^+$ were used, and with
$\begin{pmatrix}\alpha_k\\\beta^\dagger_k\end{pmatrix}
=U^\dagger_k\begin{pmatrix}c_{k\uparrow}\\c^\dagger_{-k\downarrow}\end{pmatrix}$.
Here $U_k$ is known as the Bogoliubov transformation, with the corresponding Bogoliubov
quasiparticles
\begin{align}
\alpha_k&=u_k c_{k\uparrow}-v_k c^\dagger_{-k\downarrow},\\
\beta_k&=u_k c_{-k\downarrow}+v_k c^\dagger_{k\uparrow},
\end{align}
which reproduces (5.43)--(5.47). They obey the usual fermionic anticommutation
relations, as can be seen as follows (for arbitrary unitary transformations).

\paragraph{Nambu spinors and the anticommutation relations.}
Define
\begin{equation}
\psi_k=\begin{pmatrix}c_{k\uparrow}\\c^\dagger_{-k\downarrow}\end{pmatrix},\qquad
\psi^\dagger_k=\begin{pmatrix}c^\dagger_{k\uparrow}&c_{-k\downarrow}\end{pmatrix}
\qquad\text{(Nambu spinors)},
\end{equation}
with $\{\psi_{k,i},\psi^\dagger_{k,j}\}
=\psi_{k,i}\psi^\dagger_{k,j}+\psi^\dagger_{k,j}\psi_{k,i}=\delta_{ij}$.
The Bogoliubov transformation can then be written as
\begin{equation}
\Gamma_k\equiv\begin{pmatrix}\alpha_k\\\beta^\dagger_k\end{pmatrix}=U^\dagger\psi_k,
\qquad
\Gamma^\dagger_k\equiv\psi^\dagger_k U ,
\end{equation}
so that, using bilinearity of $\{\cdot,\cdot\}$ and that the $U^{(\dagger)}_{\alpha\beta}$
are numbers,
\begin{align}
\{\Gamma_{k,i},\Gamma^\dagger_{k,j}\}
&=\Big\{\sum_\ell U^\dagger_{i,\ell}\psi_{k,\ell},\,
 \sum_m\psi^\dagger_{k,m}U_{m,j}\Big\}\\
&=\sum_{\ell,m}U^\dagger_{i,\ell}\{\psi_{k,\ell},\psi^\dagger_{k,m}\}U_{m,j}\\
&=\sum_\ell U^\dagger_{i,\ell}U_{\ell,j}=(U^\dagger U)_{ij}=\delta_{ij}.
\end{align}
Since $U$ is unitary, $\alpha_k$ and $\beta_k$ are fermionic operators.

\paragraph{Coherence factors.}
With (5.45),
\begin{align}
\big(u_kv_k\big)^2
&=\frac{1}{4}\Big(1-\frac{\varepsilon_k^2}{(E_k^+)^2}\Big)
 =\frac{1}{4}\Big(\frac{\varepsilon_k^2+|\Delta|^2-\varepsilon_k^2}
 {\varepsilon_k^2+|\Delta|^2}\Big)\\
&=\frac{1}{4}\,\frac{|\Delta|^2}{(E_k^+)^2}
 \qquad\Longrightarrow\qquad
 u_kv_k=\pm\frac{1}{2}\,\frac{|\Delta|}{E_k^+},
\end{align}
from which the rest follows, i.e.\ (5.51), where the negative solution leads to the
contradiction $1=\text{negative value}$.

\paragraph{The BCS ground state.}
Since $\alpha_k$ is fermionic,
$n^2_{\alpha,k}\equiv(\alpha^\dagger_k\alpha_k)^2=\alpha^\dagger_k\alpha_k$, and thus any
eigenvalue $\lambda$ of $n_{\alpha,k}$ must satisfy $\lambda^2=\lambda$, i.e.\
$\lambda\in\{0,1\}$ (and the same for $n_{\beta,k}$). Since $E_k^+\ge0$, it follows that
if $\alpha_k|\psi_{\mathrm{BCS}}\rangle=\beta_k|\psi_{\mathrm{BCS}}\rangle=0$, then
$|\psi_{\mathrm{BCS}}\rangle$ is the ground state of
$H_{\mathrm{eff}}=\sum_k E_k^+(\alpha^\dagger_k\alpha_k+\beta^\dagger_k\beta_k)$.

\paragraph{The pairing ansatz.}
Following Becca \& Sorella (1.91),
\begin{align}
|\psi_{\mathrm{BCS}}\rangle
&=\exp\Big(\sum_k f_k c^\dagger_{k\uparrow}c^\dagger_{-k\downarrow}\Big)|0\rangle\\
&=\prod_k\exp\big(f_k c^\dagger_{k\uparrow}c^\dagger_{-k\downarrow}\big)|0\rangle\\
&=\prod_k\Big[\sum_{n=0}^\infty\frac{f_k^n}{n!}
 \big(c^\dagger_{k,\uparrow}\big)^n\big(c^\dagger_{-k,\downarrow}\big)^n\Big]|0\rangle,
 \qquad\text{where }\big(c^\dagger_\alpha\big)^2|0\rangle=0,\\
&=\prod_k\big(1+f_k c^\dagger_{k,\uparrow}c^\dagger_{-k,\downarrow}\big)|0\rangle\\
&=\bigotimes_k\Big(|0\rangle_k
 +f_k c^\dagger_{k\uparrow}c^\dagger_{-k\downarrow}|0\rangle_k\Big),
\end{align}
with the one-momentum vacuum $|0\rangle_k$ defined by $|0\rangle=\bigotimes_k|0\rangle_k$
and $c^\dagger_{k\uparrow}|0\rangle_k=c^\dagger_{-k\downarrow}|0\rangle_k=0$. Thus this
ansatz is a pure product state, $|\psi_{\mathrm{BCS}}\rangle
=\bigotimes_k|\psi_{\mathrm{BCS}}\rangle_k$, and therefore
\begin{equation}
\alpha_k|\psi_{\mathrm{BCS}}\rangle
=\alpha_k\bigotimes_{k'}|\psi_{\mathrm{BCS}}\rangle_{k'}
=|\psi\rangle_{k_1}\otimes\dots\otimes
 \alpha_k|\psi_{\mathrm{BCS}}\rangle_k\otimes\dots
\end{equation}
so that only a single momentum sector has to be examined:
\begin{align}
\alpha_k|\psi_{\mathrm{BCS}}\rangle_k
&=\big(u_k c_{k\uparrow}-v_k c^\dagger_{-k\downarrow}\big)
 \Big(|0\rangle_k+f_k c^\dagger_{k\uparrow}c^\dagger_{-k\downarrow}|0\rangle_k\Big)\\
&=\big(u_kf_k c_{k\uparrow}c^\dagger_{k\uparrow}c^\dagger_{-k\downarrow}
 -v_k c^\dagger_{-k\downarrow}\big)|0\rangle_k\overset{!}{=}0\\
\Longleftrightarrow\quad
v_k c^\dagger_{-k\downarrow}|0\rangle_k
&=u_kf_k\big(1-c^\dagger_{k\uparrow}c_{k\uparrow}\big)
 c^\dagger_{-k\downarrow}|0\rangle_k\\
&=u_kf_k\Big(c^\dagger_{-k\downarrow}
 +c^\dagger_{k\uparrow}\underbrace{c^\dagger_{-k\downarrow}c_{k\uparrow}}_{=0}\Big)
 |0\rangle_k ,
\end{align}
where $c_\alpha|0\rangle=0$ and $c^\dagger_\alpha c^\dagger_\alpha|0\rangle=0$ were used.
Since $u_k\neq0$ this gives $f_k=v_k/u_k$, and similarly
\begin{align}
\beta_k|\psi_{\mathrm{BCS}}\rangle_k
&=\big(u_kf_k c_{-k\downarrow}c^\dagger_{k\uparrow}c^\dagger_{-k\downarrow}
 +v_k c^\dagger_{k\uparrow}\big)|0\rangle_k\overset{!}{=}0\\
\Longleftrightarrow\quad
v_k c^\dagger_{k\uparrow}|0\rangle_k
&=u_kf_k\,c^\dagger_{k\uparrow}
 \big(1-c^\dagger_{-k\downarrow}c_{-k\downarrow}\big)|0\rangle_k ,
\end{align}
again yielding $f_k=v_k/u_k$.

Therefore $|\psi_{\mathrm{BCS}}\rangle$ is indeed the ground state of the mean-field
Hamiltonian, and it behaves like the vacuum of the Bogoliubons. To avoid problems due to
$u_k=0$, one often writes instead
\begin{equation}
|\psi_{\mathrm{BCS}}\rangle
=\prod_k\big(u_k+v_k c^\dagger_{k\uparrow}c^\dagger_{-k\downarrow}\big)|0\rangle .
\end{equation}

\subsection{References}
\begin{itemize}
\item G.\ Czycholl, \emph{Solid State Theory}, Volume 2 (and Sec.~5.3 of Volume 1 for
      Kramers' theorem).
\item H.\ Bruus \& K.\ Flensberg, \emph{Many-Body Quantum Theory in Condensed Matter
      Physics}, Sec.~3.8.
\item F.\ Becca \& S.\ Sorella, \emph{Quantum Monte Carlo Approaches for Correlated
      Systems}, Eq.~(1.91).
\item M.\ Sentef, \emph{Superconductivity}, IMPRS UFAST 2017 course.
\end{itemize}

\end{document}